%%%%%
%
\section{Correction de sauts simples avec prise en compte des valeurs manquantes}

\itemstyle{Op\'erateur~:} \operfmt{correc\_saut\_mnq}
\index{correc\_saut\_mnq}

%%%%%OBJET%%%%
\itemstyle{Objet~:}

D\'etection de sauts sup\'erieurs \`a un seuil 
se produisant sur une seule valeur. La valeur "fausse"
est remplac\'ee par la moyenne des valeurs qui l'encadrent. 
D\'etection de sauts sup\'erieurs \`a un seuil\_mnq 
se produisant sur une seule valeur apr�s une valeur manquante.
 La valeur "fausse"
est remplac\'ee par une valeur manquante.

%%%%%ARGUMENTS%%%%
\bigskip
%%%%%ENTREES%%%%
\begin{tabarg}
\multicolumn{3}{l}{\titretab {Entr\'ees}}\\
\hline
$X$ & Vecteur, {\tt float} & Param\`etre d'entr\'ee \\
$S_0$ & Coeff., {\tt float} & Seuil de d\'etection d'un saut \\
$S_{mnq}$ & Coeff., {\tt float} & Seuil de d\'etection d'un saut apr�s une valeur manquante \\
\end{tabarg}

\medskip
%%%%%SORTIES%%%%
\begin{tabarg}
\multicolumn{3}{l}{\titretab {Sorties}}\\
\hline
$flag$ & Vecteur, {\tt float} & Flag de correction ou de d�tection simple\\
$X_c$ & Vecteur, {\tt float} & Param\`etre corrig\'e (m\^eme unit\'e que $X$)\\
\end{tabarg}


%%%%%FORMULE%%%%
\itemstyle{Formule}

%%correc_saut_mnq
A chaque instant $j$ :
\begin{itemize}
\item un saut est d\'etect\'e si : \newline
$X[j] - X[j-1] > S_0$ et $X[j] - X[j+1] > S_0$  (pic vers le haut)

ou 

$X[j-1] - X[j] > S_0$ et $X[j+1] - X[j] > S_0$  (pic vers le bas)

\hspace*{1cm} alors $flag[j] = 1$ et $X_c[j] = \frac{X[j+1] + X[j-1]}{2}$ \newline

\item un saut apr�s une valeur manquante est d\'etect\'e si :\newline
$X[j-1] = mnq$
et
$|X[j+1] - X[j]| > S_{mnq}$ 

\hspace*{1cm} alors $flag[j] = 2$ et $X_c[j] = mnq$ \newline

\item un saut avant une valeur manquante est d\'etect\'e si :\newline
$X[j+1] = mnq$
et
$|X[j] - X[j-1]| > S_{mnq}$ 

\hspace*{1cm} alors $flag[j] = -2$ et $X_c[j] = X[j]$ \newline

\item une forte variation positive est d\'etect\'e si :\newline
$X[j+1] - X[j] > S_0$ et $X[j] - X[j-1] > S_0$  (forte variation vers le haut)

\hspace*{1cm} alors $flag[j] = -1$ et  $X_c[j] = X[j]$ \newline

\item une forte variation n�gative est d\'etect\'e si :\newline
$X[j-1] - X[j] > S_0$ et $X[j] - X[j+1] > S_0$  (forte variation vers le bas)

\hspace*{1cm} alors $flag[j] = -1$ et $X_c[j] = X[j]$ \newline

\item aucune variation "particuli�re" n'est d�tect�e :\newline
\hspace*{1cm} alors $flag[j] = 0$ et $X_c[j] = X[j]$ \newline

\end{itemize}


%%%%%COMMENTAIRES%%%%
\itemstyle{Commentaires}

Cet op�rateur a quelques probl�mes pour : 
\begin{itemize}
\item le traitements de tous les cas de sauts avant ou apr�s une valeur manquante
\item l'interpolation lin�aire en cas de pics successifs
\item le remplacement des valeurs manquantes existantes dans $x$
\end{itemize}

Utiliser l'op�rateur \operfmt{Correc\_pics\_mnq}.


%%%%%AUTEUR%%%%
\itemstyle{Auteur~:}
CNRM/GMEI/TRAMM

%%%%%REFERENCE%%%%
\itemstyle{R\'ef\'erences~:}
